\section{第3章：动量守恒定律和能量守恒定律（3.5练习题）}
\sourceproblem{选择1}{39}{45}
\problemtext 20克小物体以500米每秒嵌入木块，碰后速度50米每秒，求木块受冲量。
\workings 小物体动量变化 $0.020(50-500)=-9\,\mathrm{N\,s}$；按相互作用冲量等大反向，木块受向前9牛秒冲量，选A。
\crosscheck 忽略其他水平冲量；碰后木块动量增加。
\sourceproblem{选择2}{40}{46}
\problemtext 质量 $m$ 的锤以向下速率 $v$ 撞桩，作用时间 $\Delta t$，重力不可忽略，求平均冲力。
\workings 若碰撞结束时锤静止，取向上为正，$\bar F\Delta t-mg\Delta t=mv$，故 $\bar F=mv/\Delta t+mg$，选A。
\crosscheck 完全非弹性碰撞本身只意味着共同末速度，并不保证静止；若末态仍向下速率 $u$，应为 $m(v-u)/\Delta t+mg$。教材选A需补充静止条件。
\sourceproblem{选择3}{40}{46}
\problemtext 质量 $m$ 的质点速率 $v$ 不变、速度方向转过 $\theta$，求净冲量大小。
\workings $I=|\Delta\bm p|=mv\sqrt{2-2\cos\theta}=2mv\sin(\theta/2)$（$0\le\theta\le\pi$），选A。
\crosscheck 方向改变即动量改变，即使动能不变。
\sourceproblem{选择4}{40}{46}
\problemtext 8千克静止物体受4牛秒净冲量，求速度大小。
\workings $v=I/m=0.5\,\mathrm{m/s}$，选A。
\crosscheck 速度方向与冲量方向一致；题目只给冲量大小。
\sourceproblem{选择5}{40}{46}
\problemtext 静止质量 $M$ 的核分出质量 $m$、速率 $v_0$ 的粒子，经典模型下求余部反冲速率。
\workings 若余部质量为 $M-m$，动量守恒得 $(M-m)v_{\rm recoil}=mv_0$，故 $v_{\rm recoil}=mv_0/(M-m)$，选A。
\crosscheck 采用教材经典质量加和近似；相对论衰变中静质量通常不守恒，需知道余部静质量并使用四动量。
\sourceproblem{选择6}{40}{46}
\problemtext 静水中船长 $L$，人与船质量相等，人由船尾走到船头，求船位移。
\workings 令船位移 $X$，人相对地面位移 $X+L$。水平质心不变，$mX+m(X+L)=0$，故 $X=-L/2$，选C。
\crosscheck 船的位移与人相对船的行走方向相反。
\sourceproblem{选择7}{40}{46}
\problemtext 静止船上两不同质量的人相向行走，判断船运动方向。
\workings 以地面系动量为准，$\bm p_{\rm boat}=-(\bm p_1+\bm p_2)$。若两人地面系动量反向且大小不等，船沿较小动量者的运动方向，选B。
\crosscheck 两人若仅给相对船的速度，需先转换参考系；若动量大小相等则船瞬时静止。
\sourceproblem{选择8}{40}{46}
\problemtext 圆锥摆在水平面内匀速率转动，判断重力和张力做功。
\workings 重力竖直，张力位于竖直径向平面，均与瞬时切向位移垂直，因此各自功为零，选A。
\crosscheck 张力水平分量虽提供向心力，但不改变动能。
\sourceproblem{选择9}{40}{46}
\problemtext 理想弹簧从自然长先拉伸5厘米，再额外拉伸5厘米，比较两段准静态外力功。
\workings $W_1=k(0.05)^2/2$，$W_2=k[(0.10)^2-(0.05)^2]/2=3W_1$，故 $W_1<W_2$，选C。
\crosscheck 若题设没有明确初始自然长，仍在初始伸长非负的线性区时后段功更大；假定准静态。
\sourceproblem{选择10}{40}{46}
\problemtext 均匀链长 $h$、质量 $m$，将下端提至上端使链对折悬挂，求外力准静态功。
\workings 初质心低于悬点 $h/2$，末质心低于悬点 $h/4$，故 $W=mg(h/2-h/4)=mgh/4$，选B。
\crosscheck 不应将整条链当作均升高 $h/2$ 的质点。
\sourceproblem{选择11}{41}{47}
\problemtext 摆长 $l$、质量 $m$，水平外力将小球缓慢从最低点移至摆角 $\theta$ 处，求功。
\workings 准静态动能不变，绳张力不做功，故外力功等于重力势能增量 $mgl(1-\cos\theta)$，选C。
\crosscheck 水平拉力通常随摆角变化，不能直接用恒力公式 $Fs$。
\sourceproblem{选择12}{41}{47}
\problemtext 比较一对作用与反作用力的冲量和功。
\workings 同一时段内 $\bm I_1=-\bm I_2$，非零时矢量不同；两力作用点位移可不同，功不必相等，选D。
\crosscheck 例：滑动接触的两物体相互摩擦功之和通常为负。零冲量为退化情形。
\sourceproblem{选择13}{41}{47}
\problemtext 弹簧原长 $l_0$，从长度 $l_1$ 伸至 $l_2$，求弹性力功。
\workings 伸长量 $x=l-l_0$，故 $W=-\int_{l_1-l_0}^{l_2-l_0}kx\,\mathrm dx=-k[(l_2-l_0)^2-(l_1-l_0)^2]/2$，选C。
\crosscheck 力与继续伸长方向相反，功为负。
\sourceproblem{选择14}{41}{47}
\problemtext 质量10千克、速度 $8\bm i+3\bm j$（m/s），求动能。
\workings $K=10(8^2+3^2)/2=365\,\mathrm J$，选C。
\crosscheck 速度分量须先平方求和。
\sourceproblem{选择15}{41}{47}
\problemtext 轮船阻力与速率平方成正比，稳定航行速率翻倍，求所需推进功率之比。
\workings 匀速时牵引力等于阻力 $kv^2$，机械功率 $P=kv^3$，翻倍后为8倍，选D。
\crosscheck 发动机输入功率还取决于推进效率；教材默认效率不变。
\sourceproblem{选择16}{41}{47}
\problemtext 质量10千克，速度 $\bm v=4t^2\bm i+16\bm j$（SI），求1至2秒合外力功。
\workings $W=\Delta K=5[(16^2+16^2)-(4^2+16^2)]=1200\,\mathrm J$，选B。
\crosscheck 亦可积分 $m\dot{\bm v}\cdot\bm v=320t^3$，结果相同。
\sourceproblem{选择17}{41}{47}
\problemtext 判断质点动量、动能、合外力冲量与功的逻辑关系。
\workings 若合外力冲量为零，$\Delta\bm p=0$，固定质量质点的 $\Delta K=0$，故合外力功为零，选C。反向不成立，因为速度方向可能变化。
\crosscheck 圆周运动半周内合力不做功却有非零冲量；必须将外力理解为合外力。
\sourceproblem{选择18}{41}{47}
\problemtext 质点从同一P点静止沿光滑固定斜面或圆弧到同一Q点，比较末动量和动能。
\workings 重力势能下降相同，约束力不做功，末速率相同；图中两条轨迹在Q处切线方向不同，故末动量矢量不同、动能相同，选D。
\crosscheck 结论依赖原图终点切线不同；若切线相同则动量也相同。
\sourceproblem{选择19}{41}{47}
\problemtext 小物体嵌入水平光滑地面上的静止木块，辨别动能损失、阻力功和生热关系（选项续PDF48）。
\workings 对入射小物体应用动能定理，其动能减少等于克服木块阻力所做的功（假定水平作用为唯一做功项），选C。系统机械能损失等于两物体总动能减少，而非入射物体单独动能减少。
\crosscheck 木块动能增加，故小物体损失的动能不全转为热，也不全转为木块动能；实际内能增加可包含形变等。
\sourceproblem{选择20}{42}{48}
\problemtext 比较圆锥摆、理想斜抛、光滑斜面自由下滑、外力牵引匀速上升四种机械能变化。
\workings 前三项在教材理想条件下机械能守恒；牵引匀速上升时动能不变而重力势能增加，外力做正功，选D。
\crosscheck 系统为物体与地球，外牵引力的功改变系统机械能。
\sourceproblem{选择21}{42}{48}
\problemtext 两物体从等高水平桌面经压缩弹簧弹出，水平射程 $s_1=s_2/2$，求质量关系。
\workings 下落时间相同，所以 $v_1=v_2/2$。弹出过程中水平动量守恒，$m_1v_1=m_2v_2$，因此 $m_1=2m_2$，选A。
\crosscheck 射程正比于初水平速率；不是正比于动能。
\sourceproblem{选择22}{42}{48}
\problemtext 竖直弹簧挂球静态伸长 $d$，由弹簧原长处静止释放，求最大伸长。
\workings $mg=kd$，最大伸长 $x$ 处速率为零，$mgx=kx^2/2$，非零根为 $x=2d$，选C。
\crosscheck 振动中心为 $d$，振幅为 $d$。
\sourceproblem{选择23}{42}{48}
\problemtext 两球完全非弹性碰撞后共同静止，比较碰前动量。
\workings 总动量为零，故 $\bm p_1=-\bm p_2$，大小相等方向相反，选D。
\crosscheck 动能为 $p^2/(2m)$，质量不同时动能不必相等。
\sourceproblem{选择24}{42}{48}
\problemtext 地球沿椭圆轨道由近日点至远日点，判断功、动能、势能及机械能变化。
\workings 距离 $r$ 增大，$U=-GMm/r$ 增大，中心引力做负功，动能减小，总机械能不变，选C。
\crosscheck 忽略其他天体扰动及非保守作用。
\sourceproblem{选择25}{42}{48}
\problemtext 孤立的质子与另一带电粒子发生库仑散射，判断系统守恒量。
\workings 教材采用静电保守相互作用、忽略辐射的两体模型，系统无外力，动量与机械能均守恒，选D。
\crosscheck 更完整的电动力学中若计辐射，须把电磁场及辐射的能量和动量纳入封闭系统；不能仅用粒子机械量声称精确守恒。
